\subsection{\textbf{RQ1 : How can Python’s dynamic nature be preserved while ensuring
  program stability through type checking?}}

Because Python is a dynamically typed language, the actual structure and state
of objects are determined only at runtime.
While this design provides significant flexibility and expressiveness for
developers, it inevitably reduces the precision attainable through static
analysis and increases the likelihood of runtime errors such as TypeError and
AttributeError.
As a result, the research community has proposed a wide range of analysis
techniques that attempt to improve program reliability while still preserving
Python’s dynamic nature.
These approaches vary widely in their assumptions, goals, and underlying
technical mechanisms—from purely static inference to hybrid and
runtime-assisted analyses—reflecting the difficulty of balancing soundness,
precision, and usability in such a dynamic environment.

\autoref{tab:RQ1} summarizes the number of existing type checking approaches,
providing an overview of the landscape. The following sections examine each
category in detail and analyze their strengthes and limitations.


\begin{table}[ht]
  \centering
  \caption{Number of papers and tools classified by Type Checking Approaches (RQ1).}
  \label{tab:RQ1}
  \small
  \setlength{\tabcolsep}{3pt} % 좌우 여백 축소
  \begin{tabularx}{\columnwidth}{
    >{\centering\arraybackslash}X|
    >{\centering\arraybackslash}X|
    >{\centering\arraybackslash}X}
    \toprule
    \textbf{Static} & \textbf{Dynamic} & \textbf{Hybrid} \\
    \midrule
    \multicolumn{1}{>{\centering\arraybackslash}X|}{21} &
    \multicolumn{1}{>{\centering\arraybackslash}X|}{2} &
    \multicolumn{1}{>{\centering\arraybackslash}X}{4} \\
    \bottomrule
  \end{tabularx}
\end{table}


\subsubsection{\textbf{Static Type Checking with Type Annotations}}


Since the introduction of type annotations, it has become possible to enhance
program reliability using explicit static type information.
Static type checkers such as Mypy~\cite{mypy} and Pyre~\cite{pyre} were
developed on this basis, enabling the detection of potential TypeError in
variables, function parameters, and return values before execution.
Moreover, Python adopts a gradual typing approach, allowing static verification
even when only parts of a program are annotated.
This enables developers to strengthen type safety incrementally.
Typically, static type checking verifies, as in
Algorithm~\autoref{alg:gradualtyping}, whether the code matches the provided
type annotations and may further perform type inference to ensure correct use
of types in practice.

\autoref{tab:RQ1} shows that static approaches constitute the largest portion
of existing research and tools.
This prevalence can be attributed to two factors.
First, static type checking imposes no runtime overhead, making it attractive
for industrial use where performance is critical.
Second, type annotation support in Python has steadily matured—particularly
with the introduction of type hints (PEP 484), generics (PEP 585), and gradual
typing infrastructures—providing a well-defined foundation upon which many
tools have been developed.
As a result, static analysis has become the most accessible and practically
deployable method for improving Python’s reliability.

However, because static analysis is limited to annotated regions, it cannot
provide the same level of safety guarantees as fully statically typed
languages. Furthermore, Python’s type annotations rely primarily on nominal
typing, which limits their ability to accurately capture Python’s structural
and behavior-based typing semantics, such as duck typing based on the presence
of specific methods or attributes. Therefore, while static type checking has
the advantage of detecting errors in advance without incurring runtime
overhead, it suffers from a high rate of false positives due to its difficulty
in capturing Python’s dynamic behavior.

% -------------------------
\begin{algorithm}[t]
\caption{GradualTypeCheck}
\begin{algorithmic}[1]
\REQUIRE Python function $f$ with optional type annotations, arguments $args$
\ENSURE TypeError if type mismatch detected
\FOR{each argument $a_i$ in $args$}
    \IF{$f$ has type annotation for $a_i$}
        \IF{type($a_i$) $\neq$ annotation\_type($a_i$)}
            \STATE raise TypeError
        \ENDIF
    \ENDIF
\ENDFOR
\STATE result = execute($f$, $args$)
\IF{$f$ has return type annotation}
    \IF{type(result) $\neq$ annotation\_type($f$.return)}
        \STATE raise TypeError
    \ENDIF
\ENDIF
\STATE \textbf{return} result
\end{algorithmic}
\label{alg:gradualtyping}
\end{algorithm}
% -------------------------


\subsubsection{\textbf{Dynamic Monitering and Runtime Enforcement}}
Dynamic type monitoring and runtime enforcement observe type-related
constraints during the actual execution of a program and report violations
immediately.
Although this approach cannot capture all errors prior to execution—as static
analysis attempts to do—it provides high accuracy for type violations along
concrete execution paths.
In languages such as Python, where dynamic binding is pervasive, this approach
is particularly practical because it can directly reproduce behavior-based
runtime errors that static techniques often fail to detect.

The most widely used runtime enforcement mechanism is decorator-based wrapping.
By attaching a type-checking decorator to a function, as in
Algorithm~\autoref{alg:monitering}, the checker verifies at each call whether
the argument and return values conform to the declared type annotations.
Libraries such as typeguard~\cite{typeguard} implement this strategy using
typing.get\_type\_hints() and isinstance-based checks.
This method requires minimal modification to existing code and detects
mis-typed values precisely at the point of failure, making it highly effective
for debugging.
However, because checks are executed on every function call, runtime overhead
is unavoidable, and complex type structures—such as generics, nested container
types—cannot always be validated at runtime.

Another form of dynamic monitoring leverages runtime traces for downstream
analyses.
MonkeyType~\cite{monkeytype}, for example, records observed argument and return
types during test execution and automatically generates Mypy-compatible
annotations.
Because this approach captures real execution behavior, it is particularly
useful for inferring types in legacy systems or analyzing complex functions.
Nevertheless, it critically depends on test quality and coverage; types for
rarely executed paths or untested branches remain unknown.

Dynamic enforcement therefore provides high precision on observed paths but
suffers from fundamental limitations.
Achieving comprehensive program safety requires sufficiently broad test or
execution coverage, creating substantial coverage dependency and potentially
high execution cost.
Furthermore, runtime information alone is often insufficient to accurately
validate deeply nested or generic types, protocol-based structural constraints,
or types involving implicit invariants.
Techniques such as conditional checks, sampling-based monitoring, or partial
instrumentation can mitigate overhead, but they cannot fully eliminate these
inherent weaknesses.

In practice, Python’s dynamic analysis ecosystem reflects these trade-offs.
typeguard offers decorator-based runtime checking, MonkeyType synthesizes
annotations from runtime traces, and CrossHair~\cite{crosshair} combines
dynamic execution with symbolic reasoning for hybrid verification. 
Collectively, these tools complement static checkers by addressing classes of
errors that remain challenging to detect in a highly dynamic language such as
Python.

% -------------------------
\begin{algorithm}[H]
\caption{RuntimeTypeMonitor}
\begin{algorithmic}[1]
\REQUIRE Python function $f$ with optional type annotations
\ENSURE Detect type violations during execution
\STATE Wrap $f$ with a type-checking decorator
\STATE On each call to $f$:
\FOR{each argument $a_i$}
    \IF{type annotation exists for $a_i$}
        \IF{type($a_i$) $\neq$ annotation\_type($a_i$)}
            \STATE Log type violation
        \ENDIF
    \ENDIF
\ENDFOR
\STATE Execute original function $f$
\IF{return type annotation exists}
    \IF{type(return) $\neq$ annotation\_type(return)}
        \STATE Log type violation
    \ENDIF
\ENDIF
\end{algorithmic}
\label{alg:monitering}
\end{algorithm}
% -------------------------


\subsubsection{\textbf{Hybrid Approaches}}
To address both Python’s dynamic nature and the inherent limitations of static
analysis, recent research has increasingly focused on hybrid analysis, which
integrates static and dynamic techniques into a unified framework.
By combining the global reasoning power of static analysis with the
execution-based precision of dynamic monitoring, hybrid methods aim to provide
more practical and robust guarantees for Python programs.

A familiar instance of such integration is the combination of gradual typing
with runtime checking.
Tools such as Mypy and Pyre statically verify annotated code regions while
validating unannotated or ambiguous segments at runtime, thereby reducing
static false positives and improving overall reliability.
Nevertheless, these approaches are not counted as hybrid in \autoref{tab:RQ1},
because they rely primarily on annotation-driven static checks combined with
simple runtime guards, rather than incorporating substantial runtime
information into the analysis process.

In \autoref{tab:RQ1}, hybrid approaches refer specifically to techniques that
use both compile-time and runtime information as essential components of the
analysis.
Examples include methods that enrich static constraints using execution traces,
that integrate dynamic attribute access patterns into static reasoning, or that
iteratively refine abstractions using runtime-observed types.
Likewise, analyses that combine symbolic reasoning with concrete execution fall
within this category.
Conversely, AI-based techniques trained solely on static code features are
excluded because they utilize only compile-time information.

By leveraging dynamic observations to refine or correct the conservatism of
static analysis, hybrid methods can reflect Python’s real-world usage patterns
more faithfully and significantly reduce spurious warnings.
At the same time, static analysis provides a global viewpoint that compensates
for the path-specific nature of dynamic monitoring, enabling broader safety
guarantees that cannot be achieved by runtime techniques alone.
This combination often yields higher type safety, improved precision, and
greater tool usability than either approach by itself.

However, hybrid methods also inherit several challenges.
Because they depend on runtime information, their effectiveness is tied to the
availability and coverage of executed paths; insufficient execution data can
limit their precision.
Runtime monitoring and trace collection introduce performance overhead, and the
integration of static and dynamic information imposes additional computational
and engineering costs. Despite these limitations, hybrid analysis remains one
of the most promising research directions for Python type analysis, offering a
balanced and practical means of handling the language’s highly dynamic
semantics.

\subsubsection{\textbf{Summary: Landscape of Type Analysis Techniques for Python}}
Because Python is fundamentally a dynamically typed language, fully static type
checking—of the kind possible in statically typed languages—is structurally
infeasible.
Object attributes may be added, removed, or modified at runtime, and even the
return type of a single function may vary depending on control flow.
Despite these challenges, recent advances in type analysis have demonstrated
that Python can achieve a meaningful degree of type safety through a
combination of complementary techniques.

As discussed in the previous sections, existing approaches fall into three
major directions, each with distinct strengths and limitations.
First, annotation-based static analysis performs ahead-of-time verification and
can detect errors early without incurring runtime costs.
Its global reasoning capability allows broad safety guarantees. However, its
inherent conservatism and difficulty in modeling Python’s dynamic behavior
often lead to false positives, particularly in unannotated or highly dynamic
regions.
Second, runtime monitoring and enforcement provide high precision for observed
execution paths and faithfully capture dynamic semantics that static methods
miss.
Yet these techniques depend heavily on execution coverage, introduce runtime
overhead, and cannot ensure global safety unless extensive tests or traces are
available.
Third, hybrid approaches integrate static and dynamic information, allowing
static analysis to be refined using concrete runtime observations.
This reduces false positives, increases practical precision, and better
reflects real-world usage patterns. At the same time, hybrid analysis still
inherits certain limitations—such as coverage dependence and the cost of
merging diverse information sources—but offers a balanced compromise between
static safety and dynamic flexibility.

In summary, while no single technique can provide complete type safety for
Python, these three categories collectively define the current landscape of
Python type analysis.
Among them, hybrid approaches are gaining prominence, as they leverage the
complementary strengths of static reasoning and dynamic evidence to address
Python’s dynamic nature more effectively than either method alone.

\begin{tcolorbox}[colback=white,colframe=blue!60,title=Answer to RQ1,fonttitle=\bfseries]
  Recent work shows a clear shift from purely static analysis toward hybrid and
  runtime-assisted approaches for Python type checking. Static analyzers still
  play a central role in large-scale verification, but runtime monitoring
  complements them by capturing behaviors that static methods miss. As a
  result, hybrid systems have become the most practical way to balance Python’s
  dynamic nature with meaningful type safety.
\end{tcolorbox}
